\subsection{积分上下限和运动分支如何选择}
上面的积分式解决三个不同问题：时间积分用来问“多久到达”，角度积分用来问“到达何处”，比耐方程用来直接求轨道形状。列式之前，先由 $E\ge U_{\rm eff}(r)$ 找允许区间，再用初始径向速度选入向或出向分支。不能因根号约定为正，就把正在靠近力心的 $\dot r$ 也取正。

\begin{origexamplebox}{从静止径向落向引力中心的时间}\label{ch11:ex:teaching-radial-fall}
质量可忽略的质点从距固定点质量 $M$ 为 $r_0$ 处静止释放，仅受万有引力。求降到 $r_1$ 的时间，其中 $0<r_1<r_0$，并讨论形式上的 $r_1\to0$ 极限。把引力源看作点质量。
\end{origexamplebox}
\begin{xesolution}
初速度为零，故角动量为零，运动是纯径向的，不可使用含 $1/L$ 的轨道公式。记 $\mu=GM$，单位质量机械能守恒给出
\[
 \tfrac12\dot r^2-\frac\mu r=-\frac\mu{r_0}.
\]
质点向内落，选负根 $\dot r=-\sqrt{2\mu(1/r-1/r_0)}$。为使时间积分为正，写成
\[
 t=\int_{r_1}^{r_0}\frac{\dd r}{\sqrt{2\mu(1/r-1/r_0)}}.
\]
令 $r=r_0\cos^2u$，$u_1=\arccos\sqrt{r_1/r_0}$。则 $|\dd r|=2r_0\sin u\cos u\,\dd u$，分母为 $\sqrt{2\mu/r_0}\tan u$，故
\[
 t=\sqrt{\frac{2r_0^3}{\mu}}\int_0^{u_1}\cos^2u\,\dd u
 =\sqrt{\frac{r_0^3}{2\mu}}\left(u_1+\sin u_1\cos u_1\right).
\]
回到半径变量，括号中的第二项为 $\sqrt{(r_1/r_0)(1-r_1/r_0)}$。当 $r_1\to r_0$，时间趋零；当 $r_1\to0$，形式极限为 $\pi\sqrt{r_0^3/(8GM)}$，量纲为时间。
现实中有半径的天体会先与质点接触；点质量模型在中心还有奇点，因此有限的形式时间不意味着质点能无条件穿过中心继续沿同一公式运动。
\end{xesolution}
